
\documentclass[11pt]{article}

\usepackage[a4paper, margin=1in]{geometry}
\usepackage[utf8]{inputenc}
\usepackage[T1]{fontenc}
\usepackage{graphicx}
\usepackage{amsmath}
\usepackage{amssymb}
\usepackage[hidelinks]{hyperref}
\usepackage{listings}
\usepackage{xcolor}
\usepackage{booktabs}
\usepackage{float}
\usepackage{longtable}
\usepackage{array}
\usepackage{titlesec}
\usepackage{fancyhdr}
\usepackage{parskip}
\usepackage{setspace}

% ============================================================
% Colors
% ============================================================

\definecolor{primary}{RGB}{30,30,120}
\definecolor{secondary}{RGB}{50,50,50}
\definecolor{accent}{RGB}{0,120,0}

% ============================================================
% Section Styling
% ============================================================

\titleformat{\section}
{\normalfont\Large\bfseries\color{primary}}
{\thesection}{1em}{}

\titleformat{\subsection}
{\normalfont\large\bfseries\color{secondary}}
{\thesubsection}{1em}{}

% ============================================================
% Header/Footer
% ============================================================

\pagestyle{fancy}
\fancyhf{}
\fancyhead[L]{Chess Game Project}
\fancyhead[R]{IBA Karachi}
\fancyfoot[C]{\thepage}

% ============================================================
% Code Listing Style
% ============================================================

\lstset{
    basicstyle=\ttfamily\small,
    keywordstyle=\color{blue},
    commentstyle=\color{green!50!black},
    stringstyle=\color{red},
    breaklines=true,
    frame=single,
    showstringspaces=false
}

% ============================================================
% Document Start
% ============================================================

\begin{document}

% ============================================================
% ============================================================
% Title Page
% ============================================================

\begin{titlepage}

\centering

% Logo
\includegraphics[width=0.5\textwidth]{IBA-logo.png}

\vspace{1cm}

% University Name
{\Large \textbf{Institute of Business Administration, Karachi}}

\vspace{0.3cm}

{\large Department of Computer Science}

\vspace{1cm}

% Project Title

\vspace{0.6cm}

{\Huge \textbf{Chess Game using C++ and SFML}}


\vspace{2cm}

% Project Details
\begin{table}[h]
\centering
\renewcommand{\arraystretch}{1.4}

\begin{tabular}{rl}
\textbf{Course:} & Object Oriented Programming \\
\textbf{Semester:} & Spring 2026 \\
\textbf{Language:} & C++17 \\
\textbf{Graphics Library:} & SFML 3.0 \\
\textbf{Instructor:} & Sir Behraj Khan \\
\end{tabular}

\end{table}

\vspace{1.5cm}

% Submitted By
{\Large \textbf{Submitted By}}

\vspace{0.5cm}

\begin{tabular}{c}
Jiya Turshani \\
Laksh Kumar \\
Deepak Raj \\
\end{tabular}

\vfill

% Bottom Section
{\large
Department of Computer Science \\
Institute of Business Administration (IBA), Karachi
}

\vspace{0.8cm}

{\large
\textbf{GitHub Repository:} \\
\url{https://github.com/deepakvalmani/Chess-main.git}
}
\vspace{0.8cm}

{\large \today}

\end{titlepage}

% ============================================================
% Table of Contents
% ============================================================

\tableofcontents
\newpage

% ============================================================
% Abstract
% ============================================================

\section{Abstract}

This project presents a complete Chess Game developed using modern C++ and the SFML graphics library. The project was developed as part of the Object Oriented Programming course at the Institute of Business Administration (IBA), Karachi.

The application implements major chess mechanics including move validation, legal movement generation, check detection, checkmate detection, stalemate detection, pawn promotion, and an interactive graphical user interface.

The project demonstrates practical implementation of Object Oriented Programming concepts including:

\begin{itemize}
    \item Encapsulation
    \item Inheritance
    \item Polymorphism
    \item Abstraction
    \item Dynamic Binding
    \item Composition
    \item Smart Pointers
\end{itemize}

The game supports both Human vs Human and Human vs AI gameplay modes.

% ============================================================
% Introduction
% ============================================================

\section{Introduction}

Chess is one of the most strategic and widely played board games in the world. Developing a chess game provides an excellent opportunity to apply advanced programming concepts and software architecture principles.

The primary goal of this project was to create a fully interactive graphical chess game using modern C++ and SFML while implementing Object Oriented Programming concepts learned during the course.

The complete system was designed from scratch including:

\begin{itemize}
    \item Chess board implementation
    \item Piece movement logic
    \item Move validation
    \item Check and checkmate detection
    \item AI opponent
    \item Graphical User Interface
\end{itemize}

% ============================================================
% Objectives
% ============================================================

\section{Objectives}

The main objectives of this project are:

\begin{enumerate}
    \item Implement a complete chess game.
    \item Apply Object Oriented Programming concepts.
    \item Design modular and reusable software architecture.
    \item Implement legal chess mechanics.
    \item Develop an AI opponent.
    \item Create a graphical interface using SFML.
\end{enumerate}

% ============================================================
% Features
% ============================================================

\section{Project Features}

\subsection{Game Modes}

\subsubsection{Play vs AI}

\begin{itemize}
    \item Human player plays as White.
    \item AI plays as Black.
    \item AI generates random legal moves.
\end{itemize}

\subsubsection{Play with Friend}

\begin{itemize}
    \item Local multiplayer support.
    \item Alternating turns between White and Black.
\end{itemize}

% ============================================================
% Chess Mechanics
% ============================================================

\section{Implemented Chess Mechanics}

\begin{table}[H]
\centering
\begin{tabular}{|l|c|}
\hline
\textbf{Feature} & \textbf{Status} \\
\hline
Pawn Movement & Implemented \\
Rook Movement & Implemented \\
Knight Movement & Implemented \\
Bishop Movement & Implemented \\
Queen Movement & Implemented \\
King Movement & Implemented \\
Capturing & Implemented \\
Move Validation & Implemented \\
Check Detection & Implemented \\
Checkmate Detection & Implemented \\
Stalemate Detection & Implemented \\
Pawn Promotion & Implemented \\
AI Opponent & Implemented \\
\hline
\end{tabular}
\caption{Implemented Chess Mechanics}
\end{table}

% ============================================================
% OOP Concepts
% ============================================================

\section{Object Oriented Programming Concepts}

\subsection{Inheritance}

All chess pieces inherit from a common base class called \texttt{Piece}.

\begin{lstlisting}[language=C++]
class Piece
{
protected:
    bool isWhite;
};
\end{lstlisting}

Derived classes include:

\begin{itemize}
    \item Pawn
    \item Rook
    \item Knight
    \item Bishop
    \item Queen
    \item King
\end{itemize}

% ============================================================

\subsection{Polymorphism}

Each chess piece overrides the movement function:

\begin{lstlisting}[language=C++]
virtual std::vector<sf::Vector2i>
getValidMoves(...) const;
\end{lstlisting}

This allows dynamic move generation depending on piece type.

% ============================================================

\subsection{Encapsulation}

Each class manages its own:

\begin{itemize}
    \item Data
    \item Movement logic
    \item Rendering logic
    \item Internal state
\end{itemize}

% ============================================================

\subsection{Composition}

The \texttt{Board} class contains all chess pieces using:

\begin{lstlisting}[language=C++]
std::unique_ptr<Piece> grid[8][8];
\end{lstlisting}

% ============================================================

\subsection{Smart Pointers}

The project uses:

\begin{lstlisting}[language=C++]
std::unique_ptr
\end{lstlisting}

Benefits include:

\begin{itemize}
    \item Automatic memory management
    \item Prevents memory leaks
    \item Safer ownership handling
\end{itemize}

% ============================================================
% Technologies
% ============================================================

\section{Technologies Used}

\begin{table}[H]
\centering
\begin{tabular}{|l|l|}
\hline
Technology & Purpose \\
\hline
C++17 & Core Programming Language \\
SFML 3.0 & Graphics and Rendering \\
CMake & Build System \\
Visual Studio & Development IDE \\
MSYS2 & Compiler Environment \\
\hline
\end{tabular}
\caption{Technologies Used}
\end{table}

% ============================================================
% Project Structure
% ============================================================

\section{Project Structure}

\begin{verbatim}
Chess/
│
├── assets/
├── include/
├── src/
├── CMakeLists.txt
├── CMakeSettings.json
└── README.md
\end{verbatim}

% ============================================================
% Setup Instructions
% ============================================================

\section{Setup Instructions}

\subsection{Requirements}

\begin{itemize}
    \item CMake 3.14+
    \item MSYS2
    \item SFML 3.0
    \item Visual Studio 2022
\end{itemize}

% ============================================================

\subsection{Installing MSYS2}

\begin{enumerate}
    \item Download MSYS2 from:
    
    \url{https://www.msys2.org/}
    
    \item Install using the UCRT64 environment.
    
    \item Run:
\end{enumerate}

\begin{lstlisting}[language=bash]
pacman -S mingw-w64-ucrt-x86_64-sfml
\end{lstlisting}

% ============================================================

\subsection{Build Steps}

\subsubsection{Step 1: Create Build Folder}

\begin{lstlisting}[language=bash]
mkdir build
cd build
\end{lstlisting}

\subsubsection{Step 2: Configure}

\begin{lstlisting}[language=bash]
cmake ..
\end{lstlisting}

\subsubsection{Step 3: Build}

\begin{lstlisting}[language=bash]
cmake --build .
\end{lstlisting}

\subsubsection{Step 4: Run}

\begin{lstlisting}[language=bash]
./Chess.exe
\end{lstlisting}


\vspace{0.5cm}

\begin{figure}[H]
        \centering
        \includegraphics[width=1.0\linewidth]{initial_screen.jpeg}
        \caption{Initial screen}
        \label{fig:placeholder}
\end{figure}


        \begin{figure}[H]
            \centering
            \includegraphics[width=1.0\linewidth]{game_screen.jpeg}
            \caption{Game Screen}
            \label{fig:placeholder}
        \end{figure}
% ============================================================
% Controls
% ============================================================

\section{Game Controls}

\begin{table}[H]
\centering
\begin{tabular}{|l|l|}
\hline
Action & Mouse Input \\
\hline
Select Piece & Left Click \\
Move Piece & Right Click \\
Return to Menu & Click Menu Button \\
\hline
\end{tabular}
\caption{Game Controls}
\end{table}

% ============================================================
% Game Logic
% ============================================================

\section{Game Logic}

\subsection{Piece Selection}

The user selects a piece using left-click.

The game:

\begin{enumerate}
    \item Detects the selected piece
    \item Generates possible moves
    \item Filters illegal moves
    \item Highlights valid moves
\end{enumerate}

% ============================================================

\subsection{Move Validation}

Each piece first generates pseudo-legal moves.

The board then verifies:

\begin{itemize}
    \item Move is inside the board
    \item Destination is valid
    \item Move does not expose the king to check
\end{itemize}

% ============================================================

\subsection{Check Detection}

The system:

\begin{enumerate}
    \item Finds the king position
    \item Generates enemy attack moves
    \item Checks whether the king is attacked
\end{enumerate}

% ============================================================

\subsection{Checkmate Detection}

A player is in checkmate when:

\begin{itemize}
    \item The king is in check
    \item No legal moves exist
\end{itemize}

% ============================================================

\subsection{AI System}

The AI:

\begin{enumerate}
    \item Collects all legal moves
    \item Stores them in a vector
    \item Randomly selects one move
    \item Executes the move
\end{enumerate}

Future improvements may include:

\begin{itemize}
    \item Minimax Algorithm
    \item Alpha-Beta Pruning
    \item Evaluation Functions
\end{itemize}

% ============================================================
% GUI Features
% ============================================================

\section{GUI Features}

The graphical interface includes:

\begin{itemize}
    \item Interactive chess board
    \item Piece textures
    \item Highlighted valid moves
    \item Status messages
    \item Turn indicators
    \item Menu system
\end{itemize}

% ============================================================
% Challenges
% ============================================================

\section{Challenges Faced}

During the development of this chess game, several technical and design challenges were encountered. These challenges helped in improving problem-solving skills and understanding of software architecture and object-oriented programming concepts.

\subsection{Implementing Chess Movement Rules}

One of the major challenges was implementing accurate movement logic for each chess piece. Every piece in chess follows different movement rules:

\begin{itemize}
    \item Pawns move forward but capture diagonally.
    \item Knights move in an ``L'' shape and can jump over pieces.
    \item Bishops move diagonally across unlimited squares.
    \item Rooks move horizontally and vertically.
    \item Queens combine rook and bishop movement.
    \item Kings move only one square in any direction.
\end{itemize}

Designing separate movement algorithms for each piece while maintaining clean object-oriented design required careful planning and testing.

\subsection{Check and Checkmate Detection}

Implementing check detection was one of the most difficult parts of the project. The system needed to determine whether the king was under attack after every move.

To solve this problem:

\begin{itemize}
    \item The board scans all enemy pieces.
    \item Each enemy piece generates its possible attack moves.
    \item The program checks whether any attack square matches the king's position.
\end{itemize}

Checkmate detection was even more complex because the game had to verify whether the player still had at least one legal move available to escape the check.

\subsection{Move Validation}

Another difficult task was preventing illegal moves. A move might appear mathematically correct but could expose the player's king to danger.

To handle this issue:

\begin{itemize}
    \item The move is temporarily simulated in memory.
    \item The board checks whether the king becomes exposed.
    \item If the king is still safe, the move is allowed.
    \item Otherwise, the move is rejected.
\end{itemize}

This required careful manipulation of piece positions using smart pointers.

\subsection{Managing Dynamic Memory}

Since chess pieces are created dynamically, memory management was an important challenge. Improper handling could cause memory leaks or crashes.

The project uses:

\begin{center}
\texttt{std::unique\_ptr}
\end{center}

to ensure automatic memory management and safe ownership transfer during captures and movement.

\subsection{SFML Rendering and GUI Design}

Creating the graphical interface using SFML also introduced challenges such as:

\begin{itemize}
    \item Correctly positioning pieces on the board
    \item Scaling textures properly
    \item Handling mouse input events
    \item Highlighting valid moves
    \item Managing menus and game states
\end{itemize}

Special attention was required to synchronize the visual board with the internal game logic.

\subsection{Debugging Complex Game Logic}

Debugging chess logic was challenging because many features depend on one another. A small issue in movement logic could affect:

\begin{itemize}
    \item Check detection
    \item Capturing
    \item AI behavior
    \item Checkmate validation
\end{itemize}

Extensive testing and debugging were required to ensure stable gameplay.

% ============================================================
% Future Improvements
% ============================================================

\section{Future Improvements}

Potential future features include:

\begin{itemize}
    \item Castling
    \item En Passant
    \item Stronger AI
    \item Sound effects
    \item Online multiplayer
    \item Move history
    \item Undo system
    \item Saving/loading games
\end{itemize}

% ============================================================
% Known Limitations
% ============================================================

\section{Known Limitations}

\begin{itemize}
    \item AI uses random move generation
    \item Castling not implemented
    \item En passant not implemented
    \item No networking support
\end{itemize}

% ============================================================
% Learning Outcomes
% ============================================================

\section{Learning Outcomes}

This project helped improve understanding of:

\begin{itemize}
    \item Object Oriented Programming
    \item Game development architecture
    \item GUI programming with SFML
    \item Event-driven programming
    \item Memory management
    \item Chess engine logic
\end{itemize}

% ============================================================
% Conclusion
% ============================================================

\section{Conclusion}

This project successfully demonstrates the implementation of a complete Chess Game using modern C++ and SFML.

The application effectively applies Object Oriented Programming concepts and demonstrates practical software engineering, GUI programming, and game logic implementation.

% ============================================================
% References
% ============================================================
\newpage
\section{References}

\begin{itemize}
    \item SFML Documentation:
    
    \url{https://www.sfml-dev.org/documentation/}
    
    \item C++ Reference:
    
    \url{https://cplusplus.com/}
    
    \item Chess Rules:
    
    \url{https://www.chess.com/learn-how-to-play-chess}
\end{itemize}

% ============================================================

\end{document}